\chapter{Living with a terminal illness}

\section*{Definition and Overview}
Living with a terminal illness is a complex medical and personal experience defined by navigating an active, progressive, and incurable disease that is expected to result in death. Terminal illnesses include advanced cancers, end-stage organ failure (heart, renal, hepatic, or respiratory), and neurodegenerative diseases such as motor neurone disease or advanced dementia. The clinical focus shifts from curative therapies to palliative and supportive care, aiming to maximise quality of life, alleviate pain and other distressing symptoms, and address the psychological, social, and spiritual needs of the individual and their family.

\section*{Epidemiology}
Globally, tens of millions of people are living with or dying from terminal illnesses annually, a number that is rising due to population ageing and medical advances that prolong survival with chronic diseases. In Australia, approximately 160,000 deaths occur each year, with the vast majority (about 70\% to 80\%) being anticipated and preceded by a period of progressive terminal decline. The most common underlying causes of terminal illness in Australia include cardiovascular disease, cancers, chronic obstructive pulmonary disease (COPD), and dementia, with dementia currently being the leading cause of death among Australian women and the second leading cause overall.

\section*{Aetiology and Pathophysiology}
The physiological progression toward the end of life varies depending on the specific disease, but it typically involves multi-organ decline and systemic metabolic changes.
\begin{itemize}
    \item \textbf{Malignant progression:} Uncontrolled cellular growth leading to local tissue destruction, metastatic spread, and systemic cachexia.
    \item \textbf{End-stage organ failure:} Gradual and irreversible loss of function in vital organs (such as the heart, lungs, kidneys, or liver), resulting in fluid overload, metabolic disturbances, or chronic hypoxia.
    \item \textbf{Neurodegeneration:} Progressive loss of neuronal structure and function, leading to cognitive decline, motor impairment, and loss of protective airway reflexes.
    \item \textbf{Systemic inflammation and cachexia:} Cytokine-mediated metabolic changes causing severe muscle wasting, anorexia, and profound weakness.
    \item \textbf{Impairment of homeostatic mechanisms:} Loss of normal physiological regulation, leading to unstable blood pressure, altered thermoregulation, and irregular breathing patterns.
\end{itemize}

\section*{Clinical Features}
The clinical features of living with a terminal illness span physical, psychological, and cognitive domains, becoming more pronounced in the final stages of life.
\begin{itemize}
    \item \textbf{Pain and physical discomfort:} Chronic nociceptive or neuropathic pain, dyspnoea, nausea, constipation, and profound fatigue.
    \item \textbf{Anorexia-cachexia syndrome:} Severe loss of appetite, weight loss, muscle wasting, and general physical frailty.
    \item \textbf{Psychological distress:} Anxiety, depression, existential distress, anticipatory grief, and fear of the dying process.
    \item \textbf{Cognitive changes:} Mild cognitive impairment, confusion, delirium (hyperactive or hypoactive), and progressive loss of consciousness.
    \item \textbf{Functional decline:} Increasing dependency on others for activities of daily living (ADLs), reduced mobility, and confinement to bed.
\end{itemize}

\section*{Differential Diagnosis}
When managing symptoms in a terminally ill patient, it is critical to distinguish between irreversible disease progression and treatable secondary conditions.
\begin{itemize}
    \item \textbf{Treatable intercurrent infections:} Conditions such as urinary tract infections or pneumonia that cause acute confusion or decline, which may be managed with conservative measures or antibiotics depending on the patient's goals of care.
    \item \textbf{Adverse medication effects:} Toxicity or side effects from drugs (such as opioid-induced delirium or sedation) that can mimic terminal disease progression.
    \item \textbf{Major depressive disorder:} Clinical depression that goes beyond normal grief or existential distress and may benefit from pharmacological or psychological intervention.
    \item \textbf{Metabolic imbalances:} Reversible acute conditions such as hypercalcaemia, dehydration, or renal impairment that can cause profound drowsiness or confusion.
    \item \textbf{Spinal cord compression:} An oncological emergency presenting with back pain and weakness that requires urgent intervention to preserve function, distinct from general weakness.
\end{itemize}

\section*{When to Seek Medical Care}
Individuals with a terminal illness require ongoing medical monitoring, with changes in management initiated as new symptoms arise or the illness progresses.
\textbf{Call 000 (or local emergency services) for:}
\begin{itemize}
    \item Sudden, severe, and distressing shortness of breath that does not respond to prescribed rescue medications.
    \item Massive, acute haemorrhage.
    \item Sudden, severe pain that cannot be controlled using the patient's breakthrough pain management plan.
    \item Acute, dangerous agitation or distress that poses a risk to the patient or carers.
\end{itemize}
Other reasons to contact the palliative care team or general practitioner include gradual worsening of pain, new-onset nausea or vomiting, changes in swallowing ability, worsening confusion, or a noticeable decline in physical mobility.

\section*{Diagnosis and Investigations}
Diagnostic investigations are carefully tailored to align with the patient's goals of care, focusing on tests that directly guide symptom control while minimising unnecessary burden.
\begin{itemize}
    \item \textbf{Comprehensive palliative assessment:} Regular, structured clinical evaluations using tools like the Palliative Care Outcomes Collaboration (PCOC) measures to assess symptom severity and carer distress.
    \item \textbf{Review of Advance Care Directives:} Formal documentation and verification of the patient's wishes, appointed decision-makers, and goals of care.
    \item \textbf{Targeted biochemistry:} Blood tests (such as calcium, renal function, or liver function) performed only if the results will change the management of symptoms like delirium or nausea.
    \item \textbf{Diagnostic imaging:} X-rays or ultrasounds used selectively, for example, to confirm pleural effusion or ascites that could be drained to relieve breathlessness or abdominal discomfort.
\end{itemize}

\section*{Management}
The management of a terminal illness is multi-disciplinary, focusing on comfort, symptom control, psychological support, and facilitating end-of-life choices.
\begin{itemize}
    \item \textbf{Palliative pharmacology:} Use of medications to control symptoms (e.g., opioids for pain and dyspnoea, antiemetics for nausea, and subcutaneous medications via syringe drivers in the final stages of life).
    \item \textbf{Advance care planning:} Documenting preferences for future medical care, including resuscitation orders and preferred place of death (home, hospice, or hospital).
    \item \textbf{Psychosocial and spiritual support:} Access to counselling, social workers, and pastoral care to help the patient and family navigate existential distress, financial concerns, and grief.
    \item \textbf{Multidisciplinary home care:} Coordination of community nursing, occupational therapy (for assistive devices), and personal care assistants to support the patient at home.
    \item \textbf{Respite care:} Temporary admission to a hospice or inpatient facility to provide carers with a break and support family coping mechanisms.
\end{itemize}

\section*{Complications}
\begin{itemize}
    \item \textbf{Refractory symptoms:} Pain, breathlessness, or delirium that does not respond to standard pharmacological management, requiring palliative sedation.
    \item \textbf{Pressure injuries:} Skin breakdown and ulceration from prolonged immobility and poor nutrition.
    \item \textbf{Carer burnout:} Physical, emotional, and psychological exhaustion experienced by family members providing round-the-clock care.
    \item \textbf{Gastrointestinal complications:} Severe constipation secondary to opioid use and immobility, or bowel obstruction in advanced abdominal malignancies.
    \item \textbf{Complicated grief:} Intense, prolonged grief reactions in surviving family members and carers.
\end{itemize}

\section*{Prognosis and Outcomes}
The prognosis is defined by the terminal nature of the underlying disease, with life expectancy varying from months to weeks or days. Prognostication is challenging and relies on clinical indicators such as functional status (e.g., the Australia-modified Karnofsky Performance Status), nutritional intake, and rate of decline. Outcomes are focused on achieving a "good death," which includes effective pain control, respecting the patient's wishes, maintaining dignity, and supporting the family through the dying process and bereavement.

\section*{Prevention and Self-Care}
\begin{itemize}
    \item Establish an advance care plan early in the course of the illness and share it with family members and healthcare providers.
    \item Take prescribed medications regularly to prevent pain and other symptoms from escalating, rather than waiting for them to become severe.
    \item Use energy conservation techniques, planning activities during periods of the day when energy levels are highest.
    \item Adapt the living environment with equipment (such as ramps, handrails, or a hospital bed) to maintain safety and independence.
    \item Stay connected with social support networks and communicate openly with loved ones about feelings and wishes.
\end{itemize}

\section*{Sources and Further Reading}
The following reputable organisations provide further information and support for individuals living with a terminal illness:
\begin{itemize}
    \item Palliative Care Australia. \textit{Resources, guidelines, and directories for palliative care services.} https://palliativecare.org.au/
    \item CareSearch. \textit{Evidence-based palliative care resources for patients, carers, and clinicians.} https://www.caresearch.com.au/
    \item Healthdirect Australia. \textit{Understanding terminal illness and palliative care.} https://www.healthdirect.gov.au/
    \item Cancer Council Australia. \textit{Living with advanced cancer support guides.} https://www.cancer.org.au/
    \item NHS. \textit{End of life care guide.} https://www.nhs.uk/
\end{itemize}
